\documentclass[12pt]{article}
\usepackage[utf8]{inputenc}
\usepackage[T1]{fontenc}
\usepackage{hyperref}
\title{Copala Triqui Language Resources for K-12 Curriculum Development}
\author{Compiled by Claude Sonnet 5 (claude-sonnet-5) for Andr\'as Kornai}
\date{}
\begin{document}
\maketitle
\section{Introduction}
Copala Triqui (ISO 639-3 trc) is an Otomanguean language of the Triqui branch, originally from the San Juan Copala area of Oaxaca, Mexico. One search result estimates 25,000 to 30,000 speakers. Diaspora communities exist, including one in the Albany, New York area.
\section{Primary Educational Resources}
No primers or school programmes were found in the search results. The online dictionary below is the most usable resource for literacy work. A paper describes it as built to help speakers become literate.
\begin{itemize}
\item Dictionary as entry to literacy and language documentation: A Copala Triqui case study: \url{https://scholarspace.manoa.hawaii.edu/items/988ecc4b-6c3b-4e78-8de9-8c7fa6dbbb3a}
\end{itemize}
\section{Linguistic and Reference Materials}
\begin{itemize}
\item Copala Triqui--Spanish--English Dictionary (Webonary), about 2,700 Triqui entries. It is credited to Roman Vidal Lopez and George Aaron Broadwell, and is free for non-commercial use with credit: \url{https://www.webonary.work/copalatriqui/?p=146817}
\item Dictionary overview and language introduction: \url{https://www.webonary.org/copalatriqui/overview/introduction/}
\item Learner's dictionary: \url{https://www.webonary.work/copalatriqui/gac987ea9-e180-470c-a541-758bc6ec5a16}
\end{itemize}
\section{Religious and Community Resources}
The paper above describes a collaboration with the diaspora organization Triquis sin Fronteras. No religious materials were found in the search.
\section{Digital and Audio Resources}
The dictionary is online and can be downloaded. No separate audio or video resources were confirmed.
\section{Curriculum Development Context}
The language is widely spoken, and a trilingual dictionary exists. No grammar texts, primers or reading materials were located in these searches. A curriculum would need to build those materials, ideally with Triqui community partners.
\end{document}